\subsection{Review design and search strategy}
The proposed design is a structured scoping review addressing the main
research question through RQ1--RQ3. It maps heterogeneous assessment concepts
and evidence gaps rather than estimating one pooled effect.
The proposal specifies English-language sources from 2017 to the final
search date, with earlier foundational sources found through citation tracing.
Confirm this scope before screening and record the actual final search date.
Candidate sources are IEEE Xplore, ACM
Digital Library, Scopus and CVF proceedings, supplemented by citation
tracing and arXiv. Identify preprints separately from peer-reviewed work.

A starting query combines (synthetic data OR synthetic images OR simulated
data) AND (computer vision OR image classification OR object detection OR
semantic segmentation) AND (quality OR evaluation OR realism OR fidelity
OR diversity OR coverage OR bias OR sim-to-real).
Companion queries will target requirements, named metrics and assessment
prerequisites so that one query does not determine the entire review.

Record each database-specific query, filters, date, result count and
export location in a separate search log. Record citation searches with
their seed papers and dates. Pilot the search against known relevant
papers before freezing it, and document later changes.
